\begin{center}\begin{tikzpicture}[font=\small,>=stealth]
\draw[black!30,rounded corners=2pt] (0,0) rectangle (14.4,8.2);
\node[anchor=west,font=\small\bfseries] at (.5,7.72) {例 3：距离和与余弦定理联立，不必分别求两边};
\begin{scope}[shift={(7.2,3.75)},scale=.58]
\draw[blue,very thick] (0,0) ellipse (5 and 4);
\draw[black!55] (-3,0)--(3,0);
\draw[orange!85!black,very thick] (-3,0)--(3.1914236925,3.0792014357) node[midway,above left] {$m$};
\draw[violet,very thick] (3,0)--(3.1914236925,3.0792014357) node[midway,left] {$n$};
\fill (-3,0) circle (2pt) node[below] {$F_1$};
\fill (3,0) circle (2pt) node[below] {$F_2$};
\fill (3.1914236925,3.0792014357) circle (2pt) node[above right] {$P$};
\node[below] at (0,0) {$6$};
\end{scope}
\node at (7.2,6.7) {$\angle F_1PF_2=60^\circ,\quad m+n=10$};

\node[anchor=west] at (.7,.55) {蓝色椭圆 $x^2/25+y^2/16=1$；三角形面积也可用底 $6$ 与高 $y_P$ 回验};
\end{tikzpicture}\end{center}
